% 第5章图：最小度与最小填充的第一步。图：K_{3,3}({3,4,6} 对 {1,2,5}) 加边{1,5}，共10条边。
% 度/缺边：1:4/3, 2:3/3, 3:3/2, 4:3/2, 5:4/3, 6:3/2。最小度选2（补{3,4},{3,6},{4,6}，最终宽度4）；最小填充选3（补{1,2},{2,5}，最终宽度3）。
\begin{tikzpicture}[x=1mm,y=1mm,>=Stealth,
  v/.style={circle,draw=FigureStroke,fill=FigureBlueFill,line width=.7pt,minimum size=5.6mm,inner sep=0pt},
  cur/.style={v,fill=FigureYellowFill,draw=FigureYellow},
  e/.style={FigureStroke,line width=.7pt},
  fe/.style={FigureGreen,line width=1pt,dashed},
  sc/.style={font=\figurefont\footnotesize,text=FigureInk},
  df/.style={font=\FigureNoteFont\fontsize{7.5}{9}\selectfont,text=FigureMuted}]
\newcommand{\drawG}[2]{%
 \node[#1] (n2) at (0,0) {$2$};
 \node[v] (n1) at (16,0) {$1$};
 \node[v] (n5) at (32,0) {$5$};
 \node[#2] (n3) at (0,18) {$3$};
 \node[v] (n4) at (16,18) {$4$};
 \node[v] (n6) at (32,18) {$6$};
 \foreach \t in {3,4,6}{\foreach \b in {1,2,5}{\draw[e] (n\t)--(n\b);}}
 \draw[e] (n1)--(n5);
 \foreach \i/\dd/\ff/\p in {1/4/3/below,2/3/3/below,5/4/3/below,3/3/2/above,4/3/2/above,6/3/2/above}{\node[df,\p=.4mm of n\i] {$\dd/\ff$};}
}
\begin{scope}
 \drawG{cur}{v}
 \draw[fe] (n3)--(n4) (n4)--(n6);
 \draw[fe] (n3) to[out=62,in=118,looseness=1.15] (n6);
 \node[sc] at (16,-17) {(a) 最小度：选$2$，补$3$条边};
 \node[df] at (16,-22) {最终宽度$4$};
\end{scope}
\begin{scope}[xshift=56mm]
 \drawG{v}{cur}
 \draw[fe] (n2)--(n1);
 \draw[fe] (n2) to[out=-62,in=-118,looseness=1.15] (n5);
 \node[sc] at (16,-17) {(b) 最小填充：选$3$，补$2$条边};
 \node[df] at (16,-22) {最终宽度$3$};
\end{scope}
\end{tikzpicture}
